\documentclass[conference]{IEEEtran}
\IEEEoverridecommandlockouts

\usepackage[spanish,es-noshorthands]{babel}
\usepackage[utf8]{inputenc}
\usepackage[T1]{fontenc}
\usepackage{amsmath,amssymb}
\usepackage{graphicx}
\usepackage{tikz}
\usepackage{pgfplots}
\pgfplotsset{compat=1.18}
\usepackage{booktabs}
\usepackage{array}
\usepackage{xcolor}
\usepackage{listings}
\usepackage{hyperref}
\usepackage{cite}
\usepackage{float}
\usepackage{caption}
\usepackage{subcaption}

\usetikzlibrary{arrows.meta, positioning, shapes.geometric, fit, backgrounds, calc}

% Colores personalizados
\definecolor{rayblue}{RGB}{30, 100, 180}
\definecolor{simdgreen}{RGB}{20, 140, 60}
\definecolor{gpuorange}{RGB}{200, 90, 20}
\definecolor{codebackground}{RGB}{248, 248, 248}
\definecolor{codekeyword}{RGB}{0, 0, 180}

% Estilo listings
\lstset{
  backgroundcolor=\color{codebackground},
  basicstyle=\ttfamily\footnotesize,
  keywordstyle=\color{codekeyword}\bfseries,
  commentstyle=\color{gray}\itshape,
  frame=single,
  breaklines=true,
  numbers=left,
  numberstyle=\tiny\color{gray},
  tabsize=2
}

% ---------------------------------------------------------------
\begin{document}

\title{\textbf{Proyecto Grupal 2: Ray Tracing Paralelo con DLP}\\
{\large Demostración 1 — CE 4302 Arquitectura de Computadores II}}

\author{
  \IEEEauthorblockN{Randall Bryan Bolaños López}
  \IEEEauthorblockA{Escuela de Ingeniería en Computadores\\
    Instituto Tecnológico de Costa Rica\\
    \textit{Semestre I, 2026}}
  \and
  \IEEEauthorblockN{Christopher Rodriguez}
  \IEEEauthorblockA{Escuela de Ingeniería en Computadores\\
    Instituto Tecnológico de Costa Rica\\
    \textit{Semestre I, 2026}}
  \and
  \IEEEauthorblockN{Andres Molina}
  \IEEEauthorblockA{Escuela de Ingeniería en Computadores\\
    Instituto Tecnológico de Costa Rica\\
    \textit{Semestre I, 2026}}
}

\maketitle

% ---------------------------------------------------------------
\begin{abstract}
Este documento presenta la propuesta del Proyecto Grupal 2 para el curso CE 4302
Arquitectura de Computadores II. Se selecciona \textit{ray tracing} como problema objetivo
por su paralelismo de datos inherente: el cómputo por píxel es independiente. Se describen
la descomposición del dominio (tiles), los diagramas para las implementaciones secuencial,
SIMD/AVX2 en CPU y CUDA en GPU, y las métricas de desempeño para su evaluación.
\end{abstract}

% ---------------------------------------------------------------
\section{Problema Seleccionado}

\subsection{Descripción}

El \textit{ray tracing} es un algoritmo de síntesis de imagen que simula el transporte de
luz mediante el trazado de rayos desde una cámara virtual hacia una escena tridimensional.
Para cada píxel $(x, y)$ de la imagen de resolución $W \times H$, se lanza al menos un rayo
primario $\mathbf{r}(t) = \mathbf{o} + t\,\mathbf{d}$, donde $\mathbf{o}$ es el origen y
$\mathbf{d}$ la dirección normalizada. El rayo se intersecta con primitivas geométricas
(esferas, triángulos) y, al detectarse una colisión, se evalúa un modelo de iluminación
(p. ej.\ Phong o Lambertian) y se emiten rayos secundarios de sombra o reflexión.

\subsection{Justificación de la Selección}

El \textit{ray tracing} es un problema \textbf{embarazosamente paralelo} (\textit{embarrassingly
parallel}) por las siguientes razones técnicas:

\begin{enumerate}
  \item \textbf{Independencia de datos por píxel}: el cálculo del color del píxel
    $(x_i, y_i)$ no depende del color calculado en $(x_j, y_j)$ para $i \neq j$.
    No existen dependencias de datos inter-píxel en la etapa de trazado primario.

  \item \textbf{Estructura uniforme de datos}: todos los rayos aplican exactamente
    las mismas operaciones aritméticas sobre distintos datos de entrada (dirección del rayo).
    Esto se mapea directamente al modelo SIMD: \textit{Single Instruction, Multiple Data}.

  \item \textbf{Alta intensidad aritmética}: las operaciones dominantes son productos
    escalares, raíces cuadradas y evaluaciones de intersección con esferas
    ($|\mathbf{d}|^2 t^2 + \ldots = 0$), todas altamente vectorizables.

  \item \textbf{Escalabilidad natural}: al incrementar la resolución
    ($1920\times1080$, $3840\times2160$) o el número de muestras por píxel
    (anti-aliasing), el volumen de trabajo crece sin cambiar la estructura de
    dependencias, permitiendo analizar escalabilidad de forma controlada.

  \item \textbf{Compatibilidad con GPU}: el modelo de ejecución SIMT
    (\textit{Single Instruction, Multiple Threads}) de CUDA mapea
    directamente un hilo por píxel, maximizando la ocupación del dispositivo.
\end{enumerate}

% ---------------------------------------------------------------
\section{Implementación SIMD en CPU (AVX2)}

\subsection{Justificación de AVX2}

El procesador base AMD Ryzen 5 3600 (6 núcleos / 12 hilos) soporta AVX2 (256 bits),
permitiendo operar sobre \textbf{8 valores \texttt{float32} simultáneamente}. No se
considera AVX-512 porque no está disponible en esta plataforma, por lo que AVX2 es el
máximo nivel SIMD explotable.

\subsection{Diagrama de Pipeline SIMD}

\begin{figure}[H]
\centering
\begin{tikzpicture}[
  simdblock/.style={rectangle, draw, fill=simdgreen!20, minimum width=5.5cm,
                    minimum height=0.6cm, align=center, font=\small},
  label/.style={font=\scriptsize\ttfamily, text=simdgreen!60!black},
  arrow/.style={-{Stealth}, thick, simdgreen!70!black}
]
  \node[simdblock] (load) {\textbf{1. Carga de 8 orígenes de rayo} (\texttt{\_mm256\_load\_ps})};
  \node[simdblock, below=0.25cm of load] (dir)
    {\textbf{2. Cálculo de 8 direcciones} — normalización vectorial};
  \node[simdblock, below=0.25cm of dir] (inter)
    {\textbf{3. Intersección rayo-esfera} (discriminante $b^2 - 4ac$)\\
     \texttt{\_mm256\_fmadd\_ps}, \texttt{\_mm256\_sqrt\_ps}};
  \node[simdblock, below=0.25cm of inter] (mask)
    {\textbf{4. Máscara de hits} — \texttt{\_mm256\_cmp\_ps}};
  \node[simdblock, below=0.25cm of mask] (shade)
    {\textbf{5. Shading condicional} — \texttt{\_mm256\_blendv\_ps}};
  \node[simdblock, below=0.25cm of shade] (store)
    {\textbf{6. Almacén al framebuffer} — \texttt{\_mm256\_store\_ps}};

  \draw[arrow] (load) -- (dir);
  \draw[arrow] (dir) -- (inter);
  \draw[arrow] (inter) -- (mask);
  \draw[arrow] (mask) -- (shade);
  \draw[arrow] (shade) -- (store);

  % Anotación
  \node[label, right=0.2cm of inter] {$\times 8$ floats};
  \node[label, right=0.2cm of mask] {Hit mask 256b};
\end{tikzpicture}
\caption{Pipeline de ejecución SIMD/AVX2: 8 rayos procesados en paralelo por ciclo de vectorización.}
\label{fig:simd}
\end{figure}

% ---------------------------------------------------------------
\section{Implementación en GPU (CUDA)}

\subsection{Mapeo de Hilos}

En CUDA se lanza un kernel con un hilo por píxel, organizado en bloques 2D de
$16 \times 16$ hilos (256 hilos/bloque), lo que favorece la coalescencia de
accesos a memoria global.

\begin{figure}[H]
\centering
\begin{tikzpicture}[scale=0.80]
  % Grid de bloques
  \foreach \bx in {0,1,2} {
    \foreach \by in {0,1} {
      \pgfmathtruncatemacro{\idx}{\bx + \by*3}
      \draw[thick, fill=gpuorange!15] ({\bx*2.2}, {\by*1.8})
        rectangle ({(\bx+1)*2.2}, {(\by+1)*1.8});
      \node[font=\scriptsize, align=center, text width=2.0cm] at ({(\bx+0.5)*2.2}, {(\by+0.5)*1.8})
        {Bloque $(\bx,\by)$\\$16{\times}16$ hilos};
    }
  }

  % Hilo individual
  \draw[thick, fill=gpuorange!50, draw=gpuorange!80] (-1.05,0.15)
    rectangle (-0.70,0.50);
  \node[font=\tiny] at (-0.875,0.325) {$(i,j)$};
  \draw[-{Stealth}, thick, gpuorange!80!black] (-0.70,0.325) -- (0.20,0.325);

  % Etiquetas
  \draw[<->, thick] (0,3.9) -- (6.6,3.9) node[midway, above] {\small $\lceil W/16 \rceil$ bloques};
  \draw[<->, thick] (7.0,0) -- (7.0,3.6) node[midway, right] {\small $\lceil H/16 \rceil$};
\end{tikzpicture}
\caption{Organización de hilos CUDA: grid 2D de bloques $16\!\times\!16$, un hilo por píxel.}
\label{fig:cuda}
\end{figure}

\subsection{Diagrama de Flujo del Kernel CUDA}

\begin{figure}[H]
\centering
\begin{tikzpicture}[
  block/.style={rectangle, draw, fill=gpuorange!20, minimum width=4.5cm,
                minimum height=0.6cm, align=center, font=\small},
  arrow/.style={-{Stealth}, thick, gpuorange!70!black}
]
  \node[block] (idx) {\textbf{Calcular $(x,y)$} desde \texttt{blockIdx}, \texttt{threadIdx}};
  \node[block, below=0.25cm of idx] (bound) {Verificar bounds $(x < W, y < H)$};
  \node[block, below=0.25cm of bound] (gen) {Generar dirección del rayo primario};
  \node[block, below=0.25cm of gen] (inter) {Loop intersección con objetos de la escena};
  \node[block, below=0.25cm of inter] (shade) {Evaluar iluminación Lambertian / Phong};
  \node[block, below=0.25cm of shade] (shadow) {Lanzar rayo de sombra (ray shadow)};
  \node[block, below=0.25cm of shadow] (write)
    {Escribir RGB al framebuffer global\\(\texttt{float3 *framebuffer[y*W+x]})};

  \draw[arrow] (idx) -- (bound);
  \draw[arrow] (bound) -- (gen);
  \draw[arrow] (gen) -- (inter);
  \draw[arrow] (inter) -- (shade);
  \draw[arrow] (shade) -- (shadow);
  \draw[arrow] (shadow) -- (write);
\end{tikzpicture}
\caption{Flujo de ejecución del kernel CUDA de ray tracing.}
\label{fig:cudaflow}
\end{figure}

% ---------------------------------------------------------------
\section{Métricas de Desempeño}

\subsection{Métricas Seleccionadas y Justificación}

La Tabla~\ref{tab:metrics} presenta las métricas elegidas junto con la
justificación técnica para el sistema objetivo.

\begin{table}[H]
\centering
\caption{Métricas de desempeño y herramientas de medición}
\label{tab:metrics}
\small
\begin{tabular}{@{}p{2.2cm}p{2.6cm}p{2.8cm}@{}}
\toprule
\textbf{Métrica} & \textbf{Herramienta} & \textbf{Justificación} \\
\midrule
Tiempo total de ejecución (ms) &
  \texttt{chrono}/CUDA Events &
  Base para speedup; comparación directa entre versiones. \\[2pt]
Speedup $S{=}T_\text{seq}/T_k$ &
  Cálculo derivado &
  Ganancia real de paralelismo vs.\ baseline secuencial. \\[2pt]
Eficiencia $E{=}S/p$ &
  Cálculo derivado &
  Fracción del potencial teórico aprovechado. \\[2pt]
Utilización SIMD (\%) &
  \texttt{perf stat}, VTune &
  Valida emisión correcta de AVX2 por el compilador. \\[2pt]
Ocupación GPU (\%) &
  Nsight Compute &
  Verifica uso de SMs; $<50\%$ requiere tuning de bloques. \\[2pt]
Ancho de banda (GB/s) &
  Nsight Systems &
  Cuello de botella potencial en resoluciones $>$1080p. \\[2pt]
Caché L1/L2 misses &
  \texttt{perf stat} &
  Guía ajuste de tamaño de tile $B_w \times B_h$. \\[2pt]
\bottomrule
\end{tabular}
\end{table}

% ---------------------------------------------------------------
\section{Conclusiones de la Propuesta}

El \textit{ray tracing} es una elección técnicamente sólida para explorar
Data-Level Parallelism: su independencia de datos por píxel lo hace ideal para
SIMD y GPU, su intensidad aritmética aprovecha las unidades vectoriales AVX2 del
Ryzen 5 3600, y la escalabilidad natural del problema permite estudiar el comportamiento
de las métricas en función del tamaño de entrada. La descomposición por dominio espacial
(tiles) provee una granularidad controlable tanto para el nivel SIMD como para el
nivel de bloques CUDA.

% ---------------------------------------------------------------
\begin{thebibliography}{9}

\bibitem{patterson2017}
J. L. Hennessy and D. A. Patterson,
\textit{Computer Architecture: A Quantitative Approach}, 6th ed.
Morgan Kaufmann, 2017.

\bibitem{shirley2020}
P. Shirley, \textit{Ray Tracing in One Weekend}, 2020. [Online].
Available: \url{https://raytracing.github.io}

\bibitem{nvidia2023}
NVIDIA Corporation, \textit{CUDA C++ Programming Guide}, v12.4, 2024.
[Online]. Available: \url{https://docs.nvidia.com/cuda/}

\bibitem{intel2023}
Intel Corporation, \textit{Intel Intrinsics Guide}, 2023.
[Online]. Available: \url{https://www.intel.com/content/www/us/en/docs/intrinsics-guide/}

\bibitem{amd2023}
AMD Inc., \textit{Software Optimization Guide},
2023.

\end{thebibliography}

\end{document}